\section{Codec Module}
\label{sec:codec}

\subsection{Purpose}

\texttt{codec.cpp} is the top-level orchestration layer.
It owns the LZKO file format, drives all other modules in the correct order, and
implements both the static and adaptive compression strategies.
The public interface (\texttt{codec.hpp}) exposes exactly two functions:
\texttt{compress()} and \texttt{decompress()}.

\subsection{Files}

\begin{itemize}
  \item \texttt{include/codec.hpp} --- public declarations.
  \item \texttt{src/codec.cpp} --- file format + all compression/decompression logic
        (\textasciitilde{}200 lines).
\end{itemize}

\subsection{Private Constants and Helpers}

\begin{lstlisting}[style=cpp]
static constexpr uint32_t BLOCK_SIZE     = 16;
static constexpr uint8_t  BLOCK_SIZE_LOG = 4;   // 2^4 = 16
\end{lstlisting}

\texttt{BLOCK\_SIZE} is defined here and not in \texttt{scanner.hpp} (see the Improvements
chapter for a discussion of this choice).

Four little-endian I/O helpers are defined as file-private (\texttt{static}) functions:

\begin{lstlisting}[style=cpp]
static void     write_u16(std::ostream&, uint16_t);
static void     write_u32(std::ostream&, uint32_t);
static uint16_t read_u16(std::istream&);
static uint32_t read_u32(std::istream&);
\end{lstlisting}

Little-endian byte order is used throughout because the primary target platform
(merlin.fit.vutbr.cz, x86-64) is natively little-endian and this matches the
internal \texttt{uint16\_t}/\texttt{uint32\_t} layout without byte swapping.

\subsubsection{pack\_block\_meta / unpack\_block\_meta}

\begin{lstlisting}[style=cpp]
static std::vector<uint8_t> pack_block_meta(const std::vector<Block>&);
static void unpack_block_meta(const std::vector<uint8_t>& buf, uint32_t n,
                               std::vector<uint8_t>& scan_dirs,
                               std::vector<uint8_t>& comp_flags);
\end{lstlisting}

These functions pack two bits per block into a byte array
(or unpack in the reverse direction).
Block $i$ occupies bit positions $2i$ (scan\_dir) and $2i+1$ (compressed),
numbered MSB-first within each byte.

Bit address formulae:
\begin{itemize}
  \item Byte index: $\lfloor 2i / 8 \rfloor = \lfloor i / 4 \rfloor$
  \item Bit position within byte (0 = MSB): $(2i) \bmod 8$
  \item Mask: \texttt{1u << (7 - bit\_position)}
\end{itemize}

\subsection{File Format}
\label{sec:file-format}

\begin{center}
\begin{tabular}{r r l}
\toprule
\textbf{Offset} & \textbf{Size} & \textbf{Field} \\
\midrule
0  & 4 B & Magic: \texttt{0x4C 0x5A 0x4B 0x4F} (``LZKO'') \\
4  & 2 B & \texttt{uint16\_t} width (little-endian) \\
6  & 2 B & \texttt{uint16\_t} height (little-endian) \\
8  & 1 B & Flags: bit~0 = \texttt{use\_model}, bit~1 = \texttt{adaptive} \\
9  & 1 B & \texttt{block\_size\_log2} (= 4) \\
10 & 4 B & \texttt{uint32\_t} original pixel count \\
14 & \ldots & Static: LZSS bit-stream \\
14 & \ldots & Adaptive: \texttt{uint32\_t num\_blocks} + metadata + per-block data \\
\bottomrule
\end{tabular}
\end{center}

Adaptive per-block data:
\begin{itemize}
  \item Packed 2-bit metadata: $\lceil 2 \cdot n / 8 \rceil$ bytes, MSB-first.
  \item Followed by $n$ entries, each: \texttt{uint32\_t stored\_size} $+$
        \texttt{stored\_size} bytes (LZSS-compressed or raw pixels).
\end{itemize}

\subsection{compress()}
\label{sec:codec-compress}

\begin{lstlisting}[style=cpp]
int compress(std::istream& in, std::ostream& out,
             int width, bool use_model, bool adaptive);
\end{lstlisting}

\textbf{Step 1: Load and validate.}
All bytes from \texttt{in} are read into \texttt{vector\textless uint8\_t\textgreater{} pixels}.
If \texttt{pixels.size() \% width != 0}, an error is printed and 1 is returned.

\textbf{Step 2: Write header.}
Magic, width, height, flags byte (\texttt{use\_model | (adaptive << 1)}),
\texttt{BLOCK\_SIZE\_LOG}, and \texttt{pixels.size()} are written.

\textbf{Step 3a: Static path} (\texttt{adaptive == false}):
\begin{enumerate}
  \item \texttt{scan\_horizontal(pixels, w, h)} $\to$ \texttt{data} (trivial copy).
  \item If \texttt{use\_model}: \texttt{forward\_model(data)}.
  \item \texttt{lzss\_encode(data, bw)}; \texttt{bw.flush()}.
\end{enumerate}

\textbf{Step 3b: Adaptive path} (\texttt{adaptive == true}), two passes:

\emph{Pass 1 --- metadata determination:}
\begin{enumerate}
  \item \texttt{split\_into\_blocks(pixels, w, h, BLOCK\_SIZE)} $\to$ \texttt{blocks}.
  \item For each block:
    \begin{itemize}
      \item Choose \texttt{scan\_dir} by comparing
            \texttt{compute\_sad\_horizontal} vs \texttt{compute\_sad\_vertical}.
      \item Scan the block pixels accordingly.
      \item If \texttt{use\_model}: \texttt{forward\_model(scanned)}.
      \item \texttt{lzss\_encode\_to\_buffer(scanned)} $\to$ \texttt{encoded}.
      \item \texttt{meta.compressed = (encoded.size() < scanned.size()) ? 1 : 0}.
    \end{itemize}
\end{enumerate}

\emph{Write metadata:} \texttt{num\_blocks} + \texttt{pack\_block\_meta(blocks)}.

\emph{Pass 2 --- data output:}
\begin{enumerate}
  \item For each block: repeat scan + model. If \texttt{meta.compressed}:
        \texttt{lzss\_encode\_to\_buffer} $\to$ write size + bytes.
        Else: write size + raw scanned bytes.
\end{enumerate}

\textbf{Two-pass rationale:}
The file format requires all block metadata to precede block data.
Because the metadata (specifically \texttt{compressed}) depends on the encoded size,
which is only known after encoding, the codec must encode each block once to determine
metadata, write all metadata, then encode again to write the data.
This doubles the LZSS encoding work for adaptive mode.
See the Improvements chapter for a fix that caches the encoded bytes from Pass 1.

\subsection{decompress()}

\begin{lstlisting}[style=cpp]
int decompress(std::istream& in, std::ostream& out);
\end{lstlisting}

\textbf{Step 1: Validate magic.}
Reads 4 bytes; if they do not match \texttt{0x4C 0x5A 0x4B 0x4F}, prints an error and
returns 1 immediately.

\textbf{Step 2: Read header.}
Extracts \texttt{width}, \texttt{height}, \texttt{flags} (from which \texttt{use\_model}
and \texttt{adaptive} are derived), \texttt{block\_size} ($= 2^{\mathtt{block\_log}}$),
and \texttt{original\_size}.

\textbf{Step 3a: Static path:}
\begin{enumerate}
  \item \texttt{lzss\_decode(BitReader(in), data, original\_size)}.
  \item If \texttt{use\_model}: \texttt{inverse\_model(data)}.
  \item Write \texttt{data} to \texttt{out}.
\end{enumerate}

\textbf{Step 3b: Adaptive path:}
\begin{enumerate}
  \item Read \texttt{num\_blocks}, then $\lceil 2 \cdot \mathtt{num\_blocks} / 8 \rceil$ metadata bytes.
  \item \texttt{unpack\_block\_meta} $\to$ \texttt{scan\_dirs[]}, \texttt{comp\_flags[]}.
  \item Allocate \texttt{image(width * height)}.
  \item For each block index \texttt{idx}:
    \begin{itemize}
      \item Compute \texttt{block\_row = idx / blocks\_x},
            \texttt{block\_col = idx \% blocks\_x}, actual \texttt{bw}, \texttt{bh}.
      \item Read \texttt{stored\_size} + raw bytes into \texttt{block\_raw}.
      \item If \texttt{comp\_flags[idx]}: wrap in \texttt{istringstream} + \texttt{BitReader}
            + \texttt{lzss\_decode(brdr, data, bw*bh)}.
            Else: \texttt{data = block\_raw}.
      \item If \texttt{use\_model}: \texttt{inverse\_model(data)}.
      \item \texttt{unscan\_horizontal} or \texttt{unscan\_vertical} based on \texttt{scan\_dirs[idx]}.
      \item Copy unscanned block pixels into the correct position in \texttt{image}.
    \end{itemize}
  \item Write \texttt{image} to \texttt{out}.
\end{enumerate}
